\section{Introducción y motivación}

La planificación de tratamientos en radioterapia moderna representa uno de los problemas computacionales más complejos dentro del área de la bioingeniería. En técnicas como la radioterapia de intensidad modulada (IMRT, \textit{Intensity Modulated Radiation Therapy}), se busca administrar una dosis alta de radiación sobre el tumor minimizando al mismo tiempo la exposición de los tejidos sanos y órganos críticos cercanos. La dificultad principal es que ambos objetivos están naturalmente en conflicto: aumentar la dosis sobre el tumor suele implicar también un incremento de la radiación absorbida por estructuras sanas.

Debido a la complejidad física y geométrica del problema, no es posible obtener soluciones analíticas exactas. Por esta razón, la planificación del tratamiento se plantea como un problema de optimización numérica de gran escala, donde las variables representan las intensidades de miles de haces elementales de radiación y las funciones objetivo permiten evaluar la calidad dosimétrica del tratamiento. El objetivo no es solamente encontrar una distribución de dosis válida, sino lograr un equilibrio adecuado entre distintos criterios clínicos que muchas veces resultan incompatibles entre sí.

Este tipo de formulaciones presenta desafíos numéricos importantes. Por un lado, la dimensión del problema es muy grande, ya que intervienen miles de parámetros y decenas de miles de puntos de evaluación espacial. Por otro, el carácter multiobjetivo del problema da lugar a múltiples soluciones posibles, conocidas en conjunto como frente de Pareto, en lugar de una única solución óptima. A esto se suman las restricciones físicas asociadas a las intensidades de radiación, que introducen dificultades adicionales durante el proceso de optimización.

Frente a este escenario, resulta necesario utilizar métodos numéricos capaces de resolver problemas de optimización multiobjetivo de gran escala con tiempos de cálculo razonables y buena estabilidad numérica. En este contexto, los métodos quasi-Newton, y en particular el algoritmo L-BFGS (\textit{Limited-memory Broyden--Fletcher--Goldfarb--Shanno}), representan una alternativa adecuada debido a su bajo costo computacional y a su reducido requerimiento de memoria.

En el presente trabajo se estudiará la formulación matemática del problema de planificación inversa en radioterapia IMRT, tomando como referencia la implementación propuesta por Lahanas et al.\ basada en el algoritmo L-BFGS. A partir de este caso de estudio se analizarán las estrategias de optimización utilizadas en problemas multiobjetivo de gran escala, el tratamiento de restricciones físicas mediante transformaciones de variables y las propiedades de convergencia y desempeño computacional del método en espacios de alta dimensionalidad.
